\section{Related literature}\label{sec:literature}

My paper connects three strands of research: the theory and evidence on minimum
wage employment effects, the evidence for developing countries with large informal
sectors, and the specific literature on Peru.

\paragraph{Theory and the canonical evidence.} The competitive model predicts that a
binding wage floor reduces employment among low-productivity workers, with an effect
whose size depends on labor demand elasticities and on the fraction of workers bound
by the floor \citep{brown_1999}. Models with employer wage-setting power deliver
different predictions: under monopsony a moderate minimum wage can raise employment,
and it typically compresses the lower tail of the wage distribution
\citep{manning_2003}. The empirical literature has moved from a presumption of sizable
disemployment toward a consensus of small effects on the margin studied. The case
study of \citet{card_krueger_1994} and the broader synthesis in \citet{card_krueger_1995}
found little evidence of job loss from the New Jersey reform, and the border-discontinuity
design of \citet{dube_lester_reich_2010} and the bunching approach of
\citet{cengiz_2019} likewise find that recent United States increases raised wages at
the bottom with limited employment effects; \citet{manning_2021} reads the accumulated
evidence as an ``elusive'' employment effect, while \citet{neumark_wascher_2007} offer
the skeptical reading. \citet{harasztosi_lindner_2019} show that a large Hungarian
increase was absorbed mostly through prices rather than employment. Closest to my
mechanism, \citet{dustmann_2022} show that the German minimum wage operated through
\emph{reallocation}---moving workers across establishments---rather than through
employment losses, and \citet{engbom_moser_2022} document how a rising Brazilian
minimum compressed the formal wage distribution; \citet{derenoncourt_montialoux_2021}
trace distributional consequences of coverage extensions in the United States.
A recurring theme, which I build on directly, is that effects should be measured
where the minimum ``bites.'' \citet{card_1992} uses regional variation in the
fraction of workers affected, \citet{clemens_wither_2019} follow workers whose wages
sat below a rising federal floor, \citet{lee_1999} uses the gap between the minimum and
the median wage, and \citet{machin_manning_rahman_2003} exploit a low-wage sector in
which the introduction of the United Kingdom minimum was especially binding. My
education-based design is an application of the same logic to a single national reform.

\paragraph{Developing countries and informality.} In economies with large informal
sectors the minimum wage interacts with the boundary between covered and uncovered
work, and two opposing forces appear repeatedly. A ``lighthouse'' effect can raise
wages even in the informal sector, where the minimum serves as a norm or reference
\citep{lemos_2009, bosch_manacorda_2010, maloney_nunez_2004}, while a reallocation
effect can push workers out of formal employment and into informal or self-employed
arrangements \citep{ham_2018_honduras}; \citet{ulyssea_2018} provides the modern
equilibrium framework in which regulation shifts firms and workers across the
formality boundary. The Latin American evidence spans this range. Early work on
Mexico and Colombia found disemployment where the floor was high relative to wages
\citep{bell_1997}, and \citet{gindling_terrell_2007} show for Costa Rica that legal
floors move wages only in the covered sector. Studies of Brazil find wage compression
with small employment effects and evidence of spillovers into the informal sector
\citep{lemos_2009, lemos_2004}; \citet{jales_2018} estimates, with a density
approach, that the Brazilian minimum moves a substantial share of workers into
informality. Work on Ecuador finds positive wage effects concentrated among
low-income workers \citep{wong_2019_ecuador}, and \citet{arango_2022_colombia}
collect the recent Colombian evidence. The Honduran case is the clearest
counterexample to the lighthouse view: \citet{ham_2018_honduras} documents that a
high minimum wage reduced covered employment and shifted workers toward informality,
with no reduction in poverty. For Uruguay, \citet{katzkowicz_2021_uruguay} use a
density-discontinuity approach in the domestic-work sector to show that effects
concentrate among the least skilled. The common methodological thread is the use of
exposure or bite measures, either geographic Kaitz indices or fraction-affected
shares, and a careful accounting of the informal margin, both of which I adopt.

\paragraph{Peru.} The domestic literature is thin and, until recently, dated; the
margin I emphasize has, however, begun to attract attention, with
\citet{corcuera_2024} studying the response of informal self-employment to Peruvian
minimum wage increases in earlier data. Early
studies of the reforms of the early 2000s find small and often statistically weak
effects. \citet{jaramillo_lopez_2006} analyze the 2003 increase using short quarterly
panels for metropolitan Lima and conclude that the reform was largely innocuous, with
negative employment-retention effects concentrated near the floor and evidence of
displacement from the formal to the informal sector. \citet{cespedes_2005} estimates a
weakly significant employment elasticity of about $-0.13$ and shows that the RMV
Granger-causes formal wages, an early statement of the lighthouse view for Peru. The
most credible modern study is \citet{jimenez_2023}, who evaluates the 2016 reform using
a difference-in-differences design with a department-level treatment dose, defined as
the pre-reform share of formal workers earning between the old and new minimum. That
study finds no disemployment, a rise in formality of 2.5 to 2.8 percentage
points, and wage gains of 4.1 to 5.4 percent, and it interprets the pattern
through employer market power, of the kind \citet{amodio_deroux_2021} document
directly in Colombian manufacturing plants. My paper differs from this precedent in three
respects that I make explicit. I study a later reform and the post-pandemic period;
I define exposure at the worker level through education rather than at the department
level; and I find a different pattern of adjustment, with no wage gain and a decline
rather than an increase in formality. I return to the reconciliation of these
findings in Section~\ref{sec:discussion}. As \citet{jimenez_2023} emphasizes, the
binding constraint for any Peruvian design is the absence of sub-national variation in
the nominal minimum, which forces the researcher to manufacture differential exposure;
my contribution is to do so through the education dimension and to trace the
consequences for the formal and informal margins.

\paragraph{Econometrics of difference-in-differences.} Methodologically I rely on the
modern difference-in-differences literature. Because my reform occurs at a single
date, the recent concerns about negative weighting and forbidden comparisons under
staggered adoption \citep{goodmanbacon_2021, dechaisemartin_2020, sun_abraham_2021,
callaway_santanna_2021} do not bind: with a common treatment date the two-way
fixed-effects estimator and its event-study analog are unbiased for a convex average
of group-time effects \citep{roth_2023_trending, baker_larcker_wang_2022}. I
nonetheless use the doubly robust estimator of \citet{santanna_zhao_2020} as a
building block, and I take the threat to identification, namely differential trends
by skill, seriously through the sensitivity analysis of \citet{rambachan_roth_2023}.
For inference with few clusters I follow \citet{bertrand_2004},
\citet{cameron_2008_bootstrap}, and \citet{pustejovsky_tipton_2018}; selection into
the wage sample is bounded with the trimming approach of \citet{lee_2009}.
Distributional effects are estimated with the unconditional quantile method of
\citet{firpo_fortin_lemieux_2009}.
